\IfFileExists{stacks-project.cls}{%
\documentclass{stacks-project}
}{%
\documentclass{amsart}
}

% For dealing with references we use the comment environment
\usepackage{verbatim}
\newenvironment{reference}{\comment}{\endcomment}
%\newenvironment{reference}{}{}
\newenvironment{slogan}{\comment}{\endcomment}
\newenvironment{history}{\comment}{\endcomment}

% For commutative diagrams we use Xy-pic
\usepackage[all]{xy}

% We use 2cell for 2-commutative diagrams.
\xyoption{2cell}
\UseAllTwocells

% We use multicol for the list of chapters between chapters
\usepackage{multicol}

% This is generally recommended for better output
\usepackage{lmodern}
\usepackage[T1]{fontenc}

% For cross-file-references
\usepackage{xr-hyper}

% Package for hypertext links:
\usepackage{hyperref}

% For any local file, say "hello.tex" you want to link to please
% use \externaldocument[hello-]{hello}
\externaldocument[introduction-]{introduction}
\externaldocument[conventions-]{conventions}
\externaldocument[sets-]{sets}
\externaldocument[categories-]{categories}
\externaldocument[topology-]{topology}
\externaldocument[sheaves-]{sheaves}
\externaldocument[sites-]{sites}
\externaldocument[stacks-]{stacks}
\externaldocument[fields-]{fields}
\externaldocument[algebra-]{algebra}
\externaldocument[brauer-]{brauer}
\externaldocument[homology-]{homology}
\externaldocument[derived-]{derived}
\externaldocument[simplicial-]{simplicial}
\externaldocument[more-algebra-]{more-algebra}
\externaldocument[smoothing-]{smoothing}
\externaldocument[modules-]{modules}
\externaldocument[sites-modules-]{sites-modules}
\externaldocument[injectives-]{injectives}
\externaldocument[cohomology-]{cohomology}
\externaldocument[sites-cohomology-]{sites-cohomology}
\externaldocument[dga-]{dga}
\externaldocument[dpa-]{dpa}
\externaldocument[sdga-]{sdga}
\externaldocument[hypercovering-]{hypercovering}
\externaldocument[schemes-]{schemes}
\externaldocument[constructions-]{constructions}
\externaldocument[properties-]{properties}
\externaldocument[morphisms-]{morphisms}
\externaldocument[coherent-]{coherent}
\externaldocument[divisors-]{divisors}
\externaldocument[limits-]{limits}
\externaldocument[varieties-]{varieties}
\externaldocument[topologies-]{topologies}
\externaldocument[descent-]{descent}
\externaldocument[perfect-]{perfect}
\externaldocument[more-morphisms-]{more-morphisms}
\externaldocument[flat-]{flat}
\externaldocument[groupoids-]{groupoids}
\externaldocument[more-groupoids-]{more-groupoids}
\externaldocument[etale-]{etale}
\externaldocument[chow-]{chow}
\externaldocument[intersection-]{intersection}
\externaldocument[pic-]{pic}
\externaldocument[weil-]{weil}
\externaldocument[adequate-]{adequate}
\externaldocument[dualizing-]{dualizing}
\externaldocument[duality-]{duality}
\externaldocument[discriminant-]{discriminant}
\externaldocument[derham-]{derham}
\externaldocument[local-cohomology-]{local-cohomology}
\externaldocument[algebraization-]{algebraization}
\externaldocument[curves-]{curves}
\externaldocument[resolve-]{resolve}
\externaldocument[models-]{models}
\externaldocument[functors-]{functors}
\externaldocument[equiv-]{equiv}
\externaldocument[pione-]{pione}
\externaldocument[etale-cohomology-]{etale-cohomology}
\externaldocument[proetale-]{proetale}
\externaldocument[relative-cycles-]{relative-cycles}
\externaldocument[more-etale-]{more-etale}
\externaldocument[trace-]{trace}
\externaldocument[crystalline-]{crystalline}
\externaldocument[spaces-]{spaces}
\externaldocument[spaces-properties-]{spaces-properties}
\externaldocument[spaces-morphisms-]{spaces-morphisms}
\externaldocument[decent-spaces-]{decent-spaces}
\externaldocument[spaces-cohomology-]{spaces-cohomology}
\externaldocument[spaces-limits-]{spaces-limits}
\externaldocument[spaces-divisors-]{spaces-divisors}
\externaldocument[spaces-over-fields-]{spaces-over-fields}
\externaldocument[spaces-topologies-]{spaces-topologies}
\externaldocument[spaces-descent-]{spaces-descent}
\externaldocument[spaces-perfect-]{spaces-perfect}
\externaldocument[spaces-more-morphisms-]{spaces-more-morphisms}
\externaldocument[spaces-flat-]{spaces-flat}
\externaldocument[spaces-groupoids-]{spaces-groupoids}
\externaldocument[spaces-more-groupoids-]{spaces-more-groupoids}
\externaldocument[bootstrap-]{bootstrap}
\externaldocument[spaces-pushouts-]{spaces-pushouts}
\externaldocument[spaces-chow-]{spaces-chow}
\externaldocument[groupoids-quotients-]{groupoids-quotients}
\externaldocument[spaces-more-cohomology-]{spaces-more-cohomology}
\externaldocument[spaces-simplicial-]{spaces-simplicial}
\externaldocument[spaces-duality-]{spaces-duality}
\externaldocument[formal-spaces-]{formal-spaces}
\externaldocument[restricted-]{restricted}
\externaldocument[spaces-resolve-]{spaces-resolve}
\externaldocument[formal-defos-]{formal-defos}
\externaldocument[defos-]{defos}
\externaldocument[cotangent-]{cotangent}
\externaldocument[examples-defos-]{examples-defos}
\externaldocument[algebraic-]{algebraic}
\externaldocument[examples-stacks-]{examples-stacks}
\externaldocument[stacks-sheaves-]{stacks-sheaves}
\externaldocument[criteria-]{criteria}
\externaldocument[artin-]{artin}
\externaldocument[quot-]{quot}
\externaldocument[stacks-properties-]{stacks-properties}
\externaldocument[stacks-morphisms-]{stacks-morphisms}
\externaldocument[stacks-limits-]{stacks-limits}
\externaldocument[stacks-cohomology-]{stacks-cohomology}
\externaldocument[stacks-perfect-]{stacks-perfect}
\externaldocument[stacks-introduction-]{stacks-introduction}
\externaldocument[stacks-more-morphisms-]{stacks-more-morphisms}
\externaldocument[stacks-geometry-]{stacks-geometry}
\externaldocument[moduli-]{moduli}
\externaldocument[moduli-curves-]{moduli-curves}
\externaldocument[examples-]{examples}
\externaldocument[exercises-]{exercises}
\externaldocument[guide-]{guide}
\externaldocument[desirables-]{desirables}
\externaldocument[coding-]{coding}
\externaldocument[obsolete-]{obsolete}
\externaldocument[fdl-]{fdl}
\externaldocument[index-]{index}

% Theorem environments.
%
\theoremstyle{plain}
\newtheorem{theorem}[subsection]{Theorem}
\newtheorem{proposition}[subsection]{Proposition}
\newtheorem{lemma}[subsection]{Lemma}

\theoremstyle{definition}
\newtheorem{definition}[subsection]{Definition}
\newtheorem{example}[subsection]{Example}
\newtheorem{exercise}[subsection]{Exercise}
\newtheorem{situation}[subsection]{Situation}

\theoremstyle{remark}
\newtheorem{remark}[subsection]{Remark}
\newtheorem{remarks}[subsection]{Remarks}

\numberwithin{equation}{subsection}

% Macros
%
\def\lim{\mathop{\mathrm{lim}}\nolimits}
\def\colim{\mathop{\mathrm{colim}}\nolimits}
\def\Spec{\mathop{\mathrm{Spec}}}
\def\Hom{\mathop{\mathrm{Hom}}\nolimits}
\def\Ext{\mathop{\mathrm{Ext}}\nolimits}
\def\SheafHom{\mathop{\mathcal{H}\!\mathit{om}}\nolimits}
\def\SheafExt{\mathop{\mathcal{E}\!\mathit{xt}}\nolimits}
\def\Sch{\mathit{Sch}}
\def\Mor{\mathop{\mathrm{Mor}}\nolimits}
\def\Ob{\mathop{\mathrm{Ob}}\nolimits}
\def\Sh{\mathop{\mathit{Sh}}\nolimits}
\def\NL{\mathop{N\!L}\nolimits}
\def\CH{\mathop{\mathrm{CH}}\nolimits}
\def\proetale{{pro\text{-}\acute{e}tale}}
\def\etale{{\acute{e}tale}}
\def\QCoh{\mathit{QCoh}}
\def\Ker{\mathop{\mathrm{Ker}}}
\def\Im{\mathop{\mathrm{Im}}}
\def\Coker{\mathop{\mathrm{Coker}}}
\def\Coim{\mathop{\mathrm{Coim}}}

% Boxtimes
%
\DeclareMathSymbol{\boxtimes}{\mathbin}{AMSa}{"02}

%
% Macros for moduli stacks/spaces
%
\def\QCohstack{\mathcal{QC}\!\mathit{oh}}
\def\Cohstack{\mathcal{C}\!\mathit{oh}}
\def\Spacesstack{\mathcal{S}\!\mathit{paces}}
\def\Quotfunctor{\mathrm{Quot}}
\def\Hilbfunctor{\mathrm{Hilb}}
\def\Curvesstack{\mathcal{C}\!\mathit{urves}}
\def\Polarizedstack{\mathcal{P}\!\mathit{olarized}}
\def\Complexesstack{\mathcal{C}\!\mathit{omplexes}}
% \Pic is the operator that assigns to X its picard group, usage \Pic(X)
% \Picardstack_{X/B} denotes the Picard stack of X over B
% \Picardfunctor_{X/B} denotes the Picard functor of X over B
\def\Pic{\mathop{\mathrm{Pic}}\nolimits}
\def\Picardstack{\mathcal{P}\!\mathit{ic}}
\def\Picardfunctor{\mathrm{Pic}}
\def\Deformationcategory{\mathcal{D}\!\mathit{ef}}
\usepackage{cleveref}
\crefname{originalsection}{original section}{original sections}
\Crefname{originalsection}{Original section}{Original sections}
\makeatletter
\newlabel{section-derived-lower-shriek-compactification}{{0F7H}{source}{}{sourcepointer-1}{}}
\newlabel{section-derived-lower-shriek-compactification@cref}{{[originalsection][0][]0F7H}{[1][0][]source}{}{}{}}
\makeatother
\crefname{originallemma}{original lemma}{original lemmas}
\Crefname{originallemma}{Original lemma}{Original lemmas}
\makeatletter
\newlabel{derived-lemma-negative-vanishing}{{05TC}{source}{}{sourcepointer-2}{}}
\newlabel{derived-lemma-negative-vanishing@cref}{{[originallemma][0][]05TC}{[1][0][]source}{}{}{}}
\makeatother
\crefname{originallemma}{original lemma}{original lemmas}
\Crefname{originallemma}{Original lemma}{Original lemmas}
\makeatletter
\newlabel{derived-lemma-unbounded-right-derived}{{07K7}{source}{}{sourcepointer-3}{}}
\newlabel{derived-lemma-unbounded-right-derived@cref}{{[originallemma][0][]07K7}{[1][0][]source}{}{}{}}
\makeatother
\crefname{originallemma}{original lemma}{original lemmas}
\Crefname{originallemma}{Original lemma}{Original lemmas}
\makeatletter
\newlabel{etale-cohomology-lemma-cohomological-dimension-proper}{{095U}{source}{}{sourcepointer-4}{}}
\newlabel{etale-cohomology-lemma-cohomological-dimension-proper@cref}{{[originallemma][0][]095U}{[1][0][]source}{}{}{}}
\makeatother
\crefname{originallemma}{original lemma}{original lemmas}
\Crefname{originallemma}{Original lemma}{Original lemmas}
\makeatletter
\newlabel{morphisms-lemma-morphism-finite-type-bounded-dimension}{{0A3V}{source}{}{sourcepointer-5}{}}
\newlabel{morphisms-lemma-morphism-finite-type-bounded-dimension@cref}{{[originallemma][0][]0A3V}{[1][0][]source}{}{}{}}
\makeatother
\title{Boundedness of the derived lower shriek functors}
\author{The Stacks Project Authors}
\date{stacks-project:more-etale.tex:local-snapshot-20261012; GNU FDL 1.2 or later}
\begin{document}
\maketitle
\setcounter{subsection}{1}
\noindent\textit{Editorial note.} All delivered mathematical text is the author's own, copied from the pinned Stacks Project chapter file more-etale.tex; no statement or proof step was rewritten, reordered or removed. The counted span is the author environment and its own proof environment. References inside the span are handed to the assembler as navigation stubs that carry the canonical tag of the cited tag only; whether each of them is load-bearing is recorded on the stub itself.\par\medskip
\noindent\textit{Editorial note.} Editorial note: in the second case (arbitrary Lambda) part (3) is argued by induction on b - a with the route fully named (truncation distinguished triangles together with Etale Cohomology tag 095U) but the induction step is not written out; part (2) is asserted to follow from part (3) with the words details omitted, the missing step being the long exact cohomology sequence bookkeeping for the truncation triangle, which is mechanical given the vanishing range established in part (3).\par\medskip
\noindent\textit{Allowed standard prerequisites (editorial).}\begin{itemize}
\item Negative vanishing for right derived functors (Derived Categories, Lemma derived-lemma-negative-vanishing, canonical tag 05TC). cited at proof lines 3816-3818. Gives the left bound in the proper case since Rf\_! = Rf\_* is a right derived functor.
\item Cohomological dimension of morphisms of finite type (Morphisms, Lemma morphisms-lemma-morphism-finite-type-bounded-dimension, canonical tag 0A3V). cited at proof lines 3819-3821. f\_* has bounded cohomological dimension for f of finite type.
\item Cohomological dimension of proper morphisms (Etale Cohomology, Lemma etale-cohomology-lemma-cohomological-dimension-proper, canonical tag 095U). cited at proof lines 3822-3823 and 3833-3834. First use: f\_* has bounded cohomological dimension for proper f, closing the torsion case. Second use: input to the induction proving part (3) in the arbitrary case.
\item Right derived functors of unbounded complexes with bounded cohomological dimension (Derived Categories, Lemma derived-lemma-unbounded-right-derived, canonical tag 07K7). cited at proof lines 3824-3825. Concludes the torsion case from the bounded cohomological dimension of f\_*.
\item Section constructing the compactification lower shriek functors (Section section-derived-lower-shriek-compactification, canonical tag 0F7H). cited at statement lines 3793-3794. The functors Rf\_! whose boundedness is claimed are the ones constructed in that section.
\end{itemize}\medskip
\ifdefined\hypertarget\hypertarget{sourcepointer-1}{}\fi
\noindent\textit{Source reference.} Original section 0F7H: sources/stacks-chow/tags.txt, source lines 19007--19007. Cited at statement lines 3793-3794: the functors Rf\_! whose boundedness is claimed are the ones constructed in that section.\par\medskip
\ifdefined\hypertarget\hypertarget{sourcepointer-2}{}\fi
\noindent\textit{Source reference.} Original lemma 05TC: sources/stacks-chow/tags.txt, source lines 7476--7476. Cited at proof lines 3816-3818: gives the left bound in the proper case since Rf\_! = Rf\_* is a right derived functor.\par\medskip
\ifdefined\hypertarget\hypertarget{sourcepointer-3}{}\fi
\noindent\textit{Source reference.} Original lemma 07K7: sources/stacks-chow/tags.txt, source lines 9647--9647. Cited at proof lines 3824-3825: concludes the torsion case from the bounded cohomological dimension of f\_*.\par\medskip
\ifdefined\hypertarget\hypertarget{sourcepointer-4}{}\fi
\noindent\textit{Source reference.} Original lemma 095U: sources/stacks-chow/tags.txt, source lines 11597--11597. Cited at proof lines 3822-3823 (torsion case: f\_* has bounded cohomological dimension for proper f) and proof lines 3833-3834 (arbitrary case: input to the induction for part (3)).\par\medskip
\ifdefined\hypertarget\hypertarget{sourcepointer-5}{}\fi
\noindent\textit{Source reference.} Original lemma 0A3V: sources/stacks-chow/tags.txt, source lines 12753--12753. Cited at proof lines 3819-3821: f\_* has bounded cohomological dimension for f of finite type.\par\medskip
\begin{lemma}
\label{lemma-derived-lower-shriek-bounded}
Let $f : X \to Y$ be a finite type separated morphism of quasi-compact
and quasi-separated schemes. Let $\Lambda$ be a ring. The functors $Rf_!$
constructed in Section \ref{section-derived-lower-shriek-compactification}
are bounded in the following sense: There exists an integer $N$ such that
for $E \in D^+_{tors}(X_\etale, \Lambda)$ or $E \in D(X_\etale, \Lambda)$
if $\Lambda$ is torsion, we have
\begin{enumerate}
\item $H^i(Rf_!(\tau_{\leq a}E) \to H^i(Rf_!(E))$ is an isomorphism
for $i \leq a$,
\item $H^i(Rf_!(E)) \to H^i(Rf_!(\tau_{\geq b - N}E))$ is an isomorphism
for $i \geq b$,
\item if $H^i(E) = 0$ for $i \not \in [a, b]$ for some
$-\infty \leq a \leq b \leq \infty$, then $H^i(Rf_!(E)) = 0$
for $i \not \in [a, b + N]$.
\end{enumerate}
\end{lemma}

\begin{proof}
Assume $\Lambda$ is torsion and consider the functor
$Rf_! : D(X_\etale, \Lambda) \to D(Y_\etale, \Lambda)$.
By construction it suffices to prove this when $f$ is an open immersion
and when $f$ is a proper morphism. For any open immersion $j : U \to X$
of schemes, the functor $j_! : D(U_\etale) \to D(X_\etale)$ is exact
and hence the statement holds with $N = 0$ in this case.
If $f$ is proper then $Rf_! = Rf_*$, i.e., it is a right derived
functor. Hence the bound on the left by
Derived Categories, Lemma \ref{derived-lemma-negative-vanishing}.
Moreover in this case $f_* : \textit{Mod}(X_\etale, \Lambda)
\to \textit{Mod}(Y_\etale, \Lambda)$ has bounded cohomological dimension by
Morphisms, Lemma \ref{morphisms-lemma-morphism-finite-type-bounded-dimension}
and \'Etale Cohomology, Lemma
\ref{etale-cohomology-lemma-cohomological-dimension-proper}.
Thus we conclude by
Derived Categories, Lemma \ref{derived-lemma-unbounded-right-derived}.

\medskip\noindent
Next, assume $\Lambda$ is arbitrary and let us consider the functor
$Rf_! : D^+_{tors}(X_\etale, \Lambda) \to D^+_{tors}(Y_\etale, \Lambda)$.
Again we immediately reduce to the case where $f$ is proper and
$Rf_! = Rf_*$. Again part (1) is immediate. To show part (3)
we can use induction on $b - a$, the distinguished
triangles of trunctions, and \'Etale Cohomology, Lemma
\ref{etale-cohomology-lemma-cohomological-dimension-proper}.
Part (2) follows from (3). Details omitted.
\end{proof}

\end{document}
